\documentclass[conference]{IEEEtran}
\IEEEoverridecommandlockouts

\usepackage{cite}
\usepackage{amsmath,amssymb,amsfonts}
\usepackage{algorithmic}
\usepackage{graphicx}
\usepackage{textcomp}
\usepackage{xcolor}
\usepackage{booktabs}
\usepackage{multirow}
\usepackage{url}
\def\BibTeX{{\rm B\kern-.05em{\sc i\kern-.025em b}\kern-.08em
    T\kern-.1667em\lower.7ex\hbox{E}\kern-.125emX}}

\begin{document}

\title{Frozen-Expert Mixture of Experts for\\Hindi-English Code-Mixed Sequence Labelling}

\author{
\IEEEauthorblockN{[Priyanshu Narayan, Samarth Agrawal, Satyasheel Singh, Vaishnav Vaijnath Pasarge]}
\IEEEauthorblockA{\textit{Deep Learning Project}\\
\textit{CS F425}\\
BITS Pilani \\
\{f20230283; f20230761; f20230261; f20230631\} :@pilani.bits-pilani.ac.in}}


\maketitle

\begin{abstract}
Hindi-English code-mixing is pervasive in Indian social media, where speakers fluidly alternate between the two languages within a single utterance. Standard monolingual models struggle with this mixed register. We present a \textbf{Frozen-Expert Mixture of Experts (FE-MoE)} framework that blends two independently pre-trained, completely frozen language models-HingBERT (Hindi-English specialist) and RoBERTa-base (English generalist)-via a lightweight learned router, training only $\sim$200K parameters. We systematically evaluate three router architectures (MLP, BiGRU, CNN) and four ablation axes (temperature, routing mode, granularity, input) across Language Identification (LID) and Part-of-Speech (POS) tagging. The best MoE variant (MLP, $\tau{=}0.1$) achieves 88.98\% weighted F1 on LID, within 0.67 points of a fully fine-tuned HingBERT (89.65\%) while using only 0.24\% as many trainable parameters. We also contribute interpretability analyses revealing that the router learns linguistically meaningful routing: Hindi tokens receive significantly higher HingBERT weight ($\alpha$) than English tokens.
\end{abstract}

\begin{IEEEkeywords}
code-mixing, mixture of experts, language identification, part-of-speech tagging, Hindi-English, frozen experts, soft routing
\end{IEEEkeywords}

% -------------------------
\section{Introduction}
% -------------------------

Hindi-English code-mixing-the practice of alternating between Hindi and English within a single sentence-is ubiquitous in Indian social media, messaging, and informal writing. Computational models for such text face a fundamental challenge: no single pre-trained language model captures both the Hindi morphology and syntax \emph{and} the English semantics equally well.

Prior work has addressed this through code-mixed pre-training \cite{nayak2022hingbert}, multilingual models, and task-specific fine-tuning. However, fine-tuning large models is compute-intensive ($\sim$85M parameters for BERT-scale models), and models fine-tuned on one task tend to forget representations useful for others.

We propose a complementary approach: \textbf{freeze both experts} and train only a small router network that learns, for each token, how much to rely on the Hindi-specialist versus the English-generalist. This draws inspiration from Soft MoE \cite{puigcerver2024softmoe} and Late Fusion in multimodal learning, but applies it at the level of independently pre-trained, task-agnostic language model encoders.

Our contributions are:
\begin{enumerate}
    \item A frozen dual-expert MoE framework with three router architectures (MLP, BiGRU, CNN) for code-mixed sequence labelling.
    \item Systematic ablations over temperature, routing mode (soft/hard/sentence-level), and router input.
    \item Interpretability analyses showing the router learns linguistically valid language-specific routing.
    \item Demonstration that 200K trainable parameters match a 85M-parameter fine-tuned model on LID.
\end{enumerate}

% -------------------------
\section{Related Work}
% -------------------------

\subsection{Mixture of Experts}

The Sparsely-Gated MoE layer \cite{shazeer2017moe} introduced the foundational design of routing tokens to expert sub-networks via a learned gating function, enabling 1000$\times$ capacity scaling. A key limitation is training instability and token dropping when routing is discrete. Puigcerver et al.\ \cite{puigcerver2024softmoe} addressed this with \emph{Soft MoE}, which replaces hard assignment with a fully-differentiable weighted combination of expert outputs-the same principle underlying our soft routing formulation. MoEBERT \cite{zuo2022moebert} adapts pre-trained BERT into an MoE structure via importance-guided distillation, achieving inference speedup while preserving knowledge. Unlike MoEBERT, we do not modify the expert internals; instead we keep both experts entirely frozen and train only the routing mechanism. MoE-LPR \cite{zhou2024moelpr} upcycles a dense LLM into MoE by freezing original FFN weights and adding expert copies with a language-prior router-closest in spirit to our method, but applied to monolingual LLM expansion rather than cross-lingual code-mixing.

\subsection{Code-Mixed NLP}

HingBERT \cite{nayak2022hingbert} is a BERT-base model pre-trained on 52.93M Hindi-English code-mixed tweets, providing strong Hindi-English representations. We use it as our Hindi-specialist expert. BERT \cite{devlin2019bert} and RoBERTa \cite{liu2019roberta} serve as the English-generalist backbone; we use RoBERTa-base as the second expert. The LinCE benchmark \cite{aguilar2020lince} established standardised evaluation for code-switching across LID, NER, and POS tasks. We evaluate on COMI-LINGUA \cite{sheth2024comilingua}, an expert-annotated large-scale dataset covering Hindi-English LID and POS with high inter-annotator agreement ($\kappa \geq 0.81$).

\subsection{Language-Aware Routing}

LR-MoE \cite{wang2023lrmoe} proposes frame-wise language routing via a shared language identification network to guide MoE expert selection in code-switched speech recognition. While our setting is text rather than speech, the principle of routing based on language identity is directly related. Our router learns to route \emph{implicitly} from hidden states rather than requiring an explicit LID signal, and we show that it recovers language-level routing without supervision on $\alpha$.

% -------------------------
\section{Methodology}
% -------------------------

\subsection{Model Architecture}

Given a sequence of $N$ words, each word $w_i$ is independently tokenized by two tokenizers: HingBERT's WordPiece and RoBERTa's BPE. Because these tokenizers produce different sub-token counts for the same word, we apply \textbf{sub-token alignment}: the hidden states of all sub-tokens belonging to the same word are averaged, yielding consistent word-level representations $\mathbf{h}^{\text{hing}}_i, \mathbf{h}^{\text{rob}}_i \in \mathbb{R}^{768}$ for each word $i$.

Both experts are kept \textbf{completely frozen} throughout training. A learned router receives these representations and produces a scalar blending weight $\alpha_i \in (0,1)$ per word:
\begin{equation}
    \hat{\mathbf{h}}_i = \alpha_i \cdot \mathbf{h}^{\text{hing}}_i + (1 - \alpha_i) \cdot \mathbf{h}^{\text{rob}}_i
\end{equation}
The blended representation $\hat{\mathbf{h}}_i$ is passed to a linear task head for per-token label prediction.

\subsection{Router Architectures}

We implement three routers, all consuming the concatenated or individual expert hidden states:

\textbf{RouterMLP} (R1--R7): Two linear layers with GELU activation, followed by sigmoid with temperature $\tau$:
\begin{equation}
    \alpha_i = \sigma\!\left(\frac{\text{Linear}(\text{GELU}(\text{Linear}([\mathbf{h}^{\text{hing}}_i; \mathbf{h}^{\text{rob}}_i]))))}{\tau}\right)
\end{equation}
Input dimensionality is $D=1536$ (concatenation) or $D=768$ (single-expert ablations).

\textbf{RouterGRU} (R8): A bidirectional GRU with hidden size 128 processes the sequence of concatenated states, capturing sequential context across positions before projecting to a scalar $\alpha$.

\textbf{RouterCNN} (R9): A 1D convolution with kernel size $k=5$ and 128 output channels captures local neighbourhood context (5-token receptive field) before projecting to $\alpha$.

\subsection{Routing Modes}

Beyond the standard token-level soft routing described above, we explore:
\begin{itemize}
    \item \textbf{Hard routing} (R7): Straight-through estimator \cite{bengio2013ste}-binary gate in the forward pass, soft gradient in the backward pass.
    \item \textbf{Sentence-level routing} (R6): $\alpha$ is computed once from the CLS tokens of both experts and broadcast to all token positions.
\end{itemize}

% -------------------------
\section{Experimental Setup}
% -------------------------

\subsection{Tasks and Datasets}

We evaluate on two tasks from the COMI-LINGUA corpus \cite{sheth2024comilingua}:
\begin{itemize}
    \item \textbf{LID}: Per-token language identification into \{hi, en, ot\}. Metric: weighted F1.
    \item \textbf{POS}: 14-class universal part-of-speech tagging. Metric: token accuracy.
\end{itemize}
Devanagari-script sentences are filtered from COMI-LINGUA to retain only romanized code-mixed text. Train/val/test splits follow the dataset's native partition with 10\% of train carved for validation (seed-controlled).

\subsection{Baselines}

\begin{itemize}
    \item \textbf{B1}: HingBERT fully fine-tuned (85M trainable params). Upper bound.
    \item \textbf{B2}: Frozen HingBERT + linear head.
    \item \textbf{B3}: Frozen RoBERTa + linear head.
    \item \textbf{B4}: Fixed 50/50 average of both experts (no router).
\end{itemize}

\subsection{Training Details}

All models use AdamW with $\text{lr}=2\times10^{-4}$, linear warmup over 10\% of steps, gradient clipping at 1.0, batch size 32, max sequence length 128, and early stopping with patience 3 (max 20 epochs). Results are reported as mean $\pm$ std over seeds \{42, 123\}.

% -------------------------
\section{Results and Ablations}
% -------------------------

\subsection{Main Results}

Table~\ref{tab:main} shows the main results. The best MoE model (R2, $\tau{=}0.1$) achieves 88.98\% LID F1, within 0.67 points of fully fine-tuned B1 (89.65\%), while training only $\sim$200K parameters versus $\sim$85M for B1-a 400$\times$ reduction. All MoE variants substantially outperform frozen single-expert baselines (B2: 80.55\%, B3: 81.48\%) and the fixed blend (B4: 80.41\%), confirming the value of learned routing.

For POS, MoE models achieve $\sim$67\% accuracy compared to B1's 91.09\%. Frozen expert representations lack the task-specific geometric structure needed for POS tagging, and the router cannot compensate without expert adaptation.

\begin{table}[t]
\centering
\caption{Main results (mean $\pm$ std, two seeds). LID: weighted F1; POS: accuracy.}
\label{tab:main}
\setlength{\tabcolsep}{4pt}
\begin{tabular}{llcc}
\toprule
\textbf{ID} & \textbf{Description} & \textbf{LID} & \textbf{POS} \\
\midrule
B1 & HingBERT fine-tuned     & \textbf{89.65} $\pm$ 0.23 & \textbf{91.09} $\pm$ 0.11 \\
B2 & HingBERT frozen          & 80.55 $\pm$ 0.06          & 65.23 $\pm$ 0.25 \\
B3 & RoBERTa frozen           & 81.48 $\pm$ 0.01          & 63.18 $\pm$ 0.31 \\
B4 & Fixed 50/50 avg          & 80.41 $\pm$ 0.31          & 61.10 $\pm$ 0.12 \\
\midrule
R1 & MoE MLP $\tau$=1.0       & 84.37 $\pm$ 3.21          & 67.76 $\pm$ 1.50 \\
R2 & MoE MLP $\tau$=0.1       & \underline{88.98} $\pm$ 0.65 & 66.99 $\pm$ 1.54 \\
R7 & MoE MLP hard routing     & 87.58 $\pm$ 0.10          & 65.20 $\pm$ 0.22 \\
R8 & MoE BiGRU router         & 86.61 $\pm$ 0.63          & 64.22 $\pm$ 2.17 \\
R9 & MoE CNN ($k$=5)          & 88.29 $\pm$ 0.12          & 66.78 $\pm$ 0.01 \\
\bottomrule
\end{tabular}
\end{table}

\subsection{Temperature Ablation}

Table~\ref{tab:ablations} (top block) shows the effect of sigmoid temperature $\tau$ on LID performance. Lower temperatures produce sharper, near-binary routing ($\tau{=}0.1$: best LID at 88.98\%) while higher temperatures collapse toward uniform blending. R1 and R5 ($\tau{=}1.0, 2.0$) exhibit high variance ($\pm$3.2\%), suggesting near-uniform routing is unstable and seed-sensitive.

\subsection{Router Architecture Ablation}

CNN (R9, 88.29\%) closely matches MLP (R2, 88.98\% at optimal $\tau$), while BiGRU (R8, 86.61\%) underperforms on LID. The CNN's local 5-token receptive field appears sufficient for code-mixing, where language alternates at the word level. The BiGRU's sequential nature may cause error propagation across language boundaries.

\subsection{Router Input Ablation}

\begin{table}[t]
\centering
\caption{Ablation results on LID (mean $\pm$ std). All MLP, $\tau$=1.0.}
\label{tab:ablations}
\setlength{\tabcolsep}{4pt}
\begin{tabular}{llcc}
\toprule
\textbf{Ablation} & \textbf{Config} & \textbf{LID} & \textbf{$\Delta$ vs R1} \\
\midrule
\multirow{2}{*}{Router input}
  & A2: HingBERT only & 86.06 $\pm$ 0.24 & $+$1.69 \\
  & A3: RoBERTa only  & 87.42 $\pm$ 0.08 & $+$3.05 \\
  & R1: Both (full)   & 84.37 $\pm$ 3.21 & -- \\
\midrule
\multirow{2}{*}{Granularity}
  & R7: Hard token    & 87.58 $\pm$ 0.10 & $+$3.21 \\
  & R6: Sentence-lvl  & 81.20 $\pm$ 0.04 & $-$3.17 \\
  & R1: Soft token    & 84.37 $\pm$ 3.21 & -- \\
\bottomrule
\end{tabular}
\end{table}

Interestingly, using only one expert's hidden states as router input (A2, A3) outperforms the full concatenation (R1) while reducing variance. RoBERTa-only routing (A3: 87.42\%) outperforms HingBERT-only (A2: 86.06\%), suggesting RoBERTa's representations carry more discriminative signal for routing even though HingBERT is the stronger expert for Hindi tokens.

Sentence-level routing (R6: 81.20\%) performs significantly worse than token-level, confirming that code-mixing requires per-token routing granularity since language identity alternates at the word level.

\subsection{Interpretability: What Does the Router Learn?}

Analysis of $\alpha$ distributions on the LID test set reveals clear language-level routing. Hindi tokens (\texttt{hi}) receive significantly higher $\alpha$ (higher HingBERT weight) than English tokens (\texttt{en}), with \texttt{ot} (Other) tokens falling between. This structure emerges without any direct supervision on $\alpha$-the router learns language-specific routing purely from the task loss. Router $\alpha$ values also correlate with expert hidden-state disagreement: tokens where $\|\mathbf{h}^{\text{hing}} - \mathbf{h}^{\text{rob}}\|_2$ is large are routed more decisively (sharper $\alpha$), while tokens both models agree on show near-uniform blending.

% -------------------------
\section{Failure Cases and Limitations}
% -------------------------

\textbf{POS tagging}: Frozen expert representations are not adapted for POS, resulting in $\sim$67\% accuracy versus B1's 91\%. The router cannot compensate for inadequate task-specific geometry in the frozen hidden states. Partially unfreezing the top transformer layers would likely close this gap.

\textbf{High variance at $\tau{=}1.0$}: R1 and R5 show $\pm$3.2\% variance across seeds, indicating the routing landscape at moderate temperatures is non-convex and sensitive to initialisation. Lower temperatures (R2, $\tau{=}0.1$) stabilize training significantly ($\pm$0.65\%).

\textbf{NER}: All frozen-expert models achieve near-zero entity-level F1, consistent with the finding that boundary detection requires expert adaptation. NER is excluded from the main sweep.

\textbf{Compute}: B1 checkpoints (418 MB) cannot be version-controlled in standard git repositories due to size constraints, limiting reproducibility without LFS or cloud storage.

% -------------------------
\section{Conclusion}
% -------------------------

We presented FE-MoE, a frozen dual-expert MoE framework for Hindi-English code-mixed sequence labelling. By keeping two complementary pre-trained experts frozen and training only a lightweight router, we achieve 88.98\% weighted F1 on LID (within 0.67 points of full fine-tuning) using just 200K trainable parameters. Across three router architectures and four ablation axes, we find that: (i) lower temperature sharpens routing and improves both performance and stability; (ii) CNN-based routing matches MLP at lower sequential cost; (iii) token-level granularity is essential for code-mixing; and (iv) the router learns implicit language-aware routing without explicit LID supervision. The FE-MoE framework offers a practical, parameter-efficient alternative to full fine-tuning for code-mixed NLP, and the interpretability analyses provide insight into how language specialisation emerges from task-driven training.

% -------------------------
\bibliographystyle{IEEEtran}
\bibliography{references}

\end{document}
